\begin{thm1}
\makeatletter\def\@currentlabel{A}\makeatother
\label{thm: desc basis are the same} 
    The set 
    \begin{align}\label{eq: our basis}
        \mathbf{x}^{\mathcal{C}_\mu} = \{\mathbf{x}^\alpha \ | \ \alpha \in \mathcal{C}_\mu\}
    \end{align}
     is a monomial basis of $R_{\mu}$. 
       In fact, it coincides with the basis $  \mathbf{x}^{\mathcal{D}_{\mu^t}}$ given by Carlsson--Chou \cite{CC}, i.e. $\mathcal{C}_\mu = \mathcal{D}_{\mu^t}$.
\end{thm1}

Before we prove Theorem \ref{thm: desc basis are the same}, we point out the connections between this basis and the Hilbert series $\Hilbq(R_{\mu})$.
Using combinatorial formula \eqref{eq:catabolizable HL formula} for the modified Hall-Littlewood polynomials and Proposition \ref{thm: charge cocharge trabspose} gives us the following expression for $\mHL_{\mu}[X;q]$.
\begin{align}
    \mHL_{\mu} [X;q] &=  \sum\limits_{\substack{T\in \syt_n\\  \ctype(T^t)\unrhd\mu}}  q^{\charge(T)} s_{\shape(T^t)}(X).
\end{align}

From this, we can see that the classification of insertion tableaux that appears in \eqref{def:C_mu} is the natural one to consider when looking at subsets of descent monomials. 
In fact, it is evident from our construction that the degrees of the monomials in our set match what we expect from $\Hilbq(R_{\mu})$. 
Rephrasing \eqref{eq:Hilb of R_mu} for $\mHL_{\mu}[X;q]$ using Proposition \ref{thm: charge cocharge trabspose}, we have
\begin{align}
     \Hilbq(R_{\mu}) &=  \sum\limits_{\substack{w\in S_n\\  \ctype(P(w)^t)\unrhd\mu}}  q^{\text{charge}(w)} .
\end{align}
It is also apparent that $|\mathcal{C}_\mu| = \dim(R_{\mu})$, which was not obvious from the definition of $\mathcal{D}_{\mu^t}$.
Furthermore, it is clear that our set has the correct number of monomials of each degree to be a basis of $R_{\mu}$ by construction.
That is, for any nonnegative $d$, we have 
\[|\{w\in S_n, \ctype(P(w))^t \unrhd \mu, \charge(w) = d\}| = \dim((R_{\mu})_d),\]
where $(R_{\mu})_d$ is the degree $d$ component of $R_{\mu}$.
This construction also gives us a direct way to determine for which $w\in S_n$ we have $g_w(\mathbf{x})\in \mathbf{x}^{\mathcal{C}_\mu}$.
In particular, we have $g_w(\mathbf{x})\in \mathbf{x}^{\mathcal{C}_\mu}$ for $w = \text{rev}(\sigma^{-1})$ where $\ctype(P(\sigma)^t)\unrhd \mu$.


\begin{ex}
Consider $\mu = (2,1,1)$. There are five standard tableaux $S$ such that $\ctype(S^t) \unrhd (2,1,1) = \mu$. We list them, along with all words $w$ such that $P(w) = S$ and their charge monomials:
\begin{center}
     \ytableausetup{smalltableaux,  aligntableaux = center}
   \begin{tabular}{ p{1cm}|p{5cm}  p{5cm}  }
$S$ & $\{w \ | \ P(w) = S\}$ & $\{x^{\cw(w)} \ | \ P(w) = S\}$\\
 \hline
\\[-.5em]
 \ytableaushort{2,134}  &\{2134,
  \ 2314, \ 2341\}& $\{x_3x_4^2,
  \ x_2x_4^2, \ x_2x_3^2\}$ \\[-.5em] \\
\ytableaushort{24,13} &  \{2143 ,\  2413\} &  $\{x_3x_4 ,\  x_2x_4\}$  \\[-.5em] \\
\ytableaushort{4,2,13} & \{4213, \ 4231, \ 2431\} &  $\{x_1x_4, \ x_1x_3, \ x_2x_3\}$\\[-.5em] \\
\ytableaushort{3,2,14} & \{3214, \ 3241 , \ 3421\} & $\{x_4, \ x_3 , \ x_2\}$\\[-.5em] \\
\ytableaushort{4,3,2,1}&  \{4321\}& $\{1\}$
\end{tabular}
\end{center}
Note that all charge monomials for words with the same insertion tableau $S$ have degree equal to charge$(S)$. 
The resulting set of charge monomials is  \[\{x_3x_4^2,
  \ x_2x_4^2, \ x_2x_3^2, \ x_3x_4 ,\  x_2x_4, \ x_1x_4, \ x_1x_3, \ x_2x_3, \ x_4, \ x_3 , \ x_2, \ 1\},\]
which is the same as in Example \ref{ex: 31}.
\end{ex}

Now we will prove Theorem \ref{thm: desc basis are the same} by showing $\mathcal{C_\mu} = \mathcal{D}_{\mu^t}$, which implies that our set is a basis of $R_{\mu}$ by the work of Carlsson--Chou.
Since we know that $\mathcal{D}_{\mu^t}$ is a basis of $R_{\mu}$ and that $\mathcal{C_\mu}$ has the correct cardinality to be a basis, it suffices to show that $\mathcal{C}_\mu\subset \mathcal{D}_{\mu^t}$.
We first make the following observation, which is a corollary of Lemma \ref{lemma:charge and majt} using \eqref{eq: charge is reverse of cocharge of reverse}:

\begin{coro}\label{cor:cocharge and majt} $D_n = \{\rev(\cc(w)) \ | \ w \in S_n\}.$
\end{coro}

To show that $\mathcal{C}_\mu\subset \mathcal{D}_{\mu^t}$, we translate the problem into one about cocharge words to make use of Algorithm \ref{alg: cat insertion}. We first consider the case where $\ctype(w) = \mu$.

\begin{prop}\label{prop: ctype implies shuffle}
    Let $\mu$ be a partition of $n$ of length $l$ and  $w$ be a permutation of $n$ with $\ctype(w) = \mu$. 
    Then $\cc(w) \in \sh'(z^{(1)},\dots , z^{(l)})$ for some $(z^{(1)},\dots , z^{(l)})\in D_{\mu^t_1}\times \cdots \times D_{\mu^t_l}$.
\end{prop}

\begin{proof}
    We obtain a filling $T_w$ of shape $\ctype(w) = \mu $ by applying Algorithm \ref{alg: cat insertion} to $\cc(w)$.
    It is clear that $\cc(w) \in \sh (\cc(w)^{(1)}, \cc(w)^{(2)}, \dots, \cc(w)^{(l)})$, where $\cc(w)^{(j)}$ is the subword of $\cc(w)$ corresponding to column $j$ in $T_w$. We know $\cc(W)^{(j)}$ has length $\mu_j$, which is the height of column $j$. 
    Thus, from Proposition \ref{prop: subword of column is a cocharge word}, we know $\cc(w)^{(j)} = \cc(\sigma)$ for $\sigma\in S_{\mu^t_j}$. 
    By Corollary \ref{cor:cocharge and majt}, this implies  $\rev(\cc(w)^{(j)}) \in \mathcal{D}_{\mu^t_j}$.
\end{proof}

\begin{ex}\label{ex: cat ex}
    Recall that in Example \ref{ex: cat alg}, we showed that for $w = 6 \ 3 \ 4  \ 1 \ 2\ 5 $, with $\cc(w) = 2 \ 1 \ 1  \ 0 \ 0\ 1 $, we have  $\ctype(w) = (2,2,2).$ and that the filling $T_w$ of $ (2,2,2) $ is 
        \[\ytableaushort{16,32,54}. \] 
    From this, we construct two subwords of lengths $3$ respectively:
    \begin{align*}
        \cc(w)^{(1)} &= \cc(w)_1\cc(w)_3\cc(w)_5 = 210\\
         \cc(w)^{(2)} &= \cc(w)_2\cc(w)_4\cc(w)_6 = 101
    \end{align*}
    We can see that $\cc(w)^{(1)} = \cc(321) = \rev(\cw(123))$ and $\cc(w)^{(2)} = \cc(213) = \rev(\cw(213))$. 
    Thus $\cc(w)\in \text{Sh} '(\cw(123), \cw(213))$ , where $(\cw(123), \cw(213)) \in \mathcal{D}_{(3,3)}$.
\end{ex}

We rephrase Proposition \ref{prop: ctype implies shuffle} to match the language of Theorem \ref{thm: desc basis are the same} using \eqref{eq: charge is reverse of cocharge of reverse}:

\begin{coro}\label{cor:ctype implies shuffle}
    For any $w\in S_n$ such that $\ctype(P(w)^t) =\mu$, we have $\cw(w)\in \mathcal{D}_{\mu^t}$.
\end{coro}
\begin{proof}
    From  \eqref{eq: charge is reverse of cocharge of reverse} we know  \[\cc(w) \in \sh'(z^{(1)},\dots , z^{(l)}) \Leftrightarrow\cw(\rev(w))= \rev(\cc(w)) \in \sh(z^{(1)},\dots , z^{(l)}).\]
    We also have $\ctype(\rev(w)) = \ctype(P(w)^t)$ using Theorem \ref{thm:RSK rev}.
\end{proof}

Note that Corollary \ref{cor:ctype implies shuffle} only considers permutations $w$ with $\ctype(P(w))^t) = \mu$. To prove the inclusion $\mathcal{C}_\mu \subset \mathcal{D}_{\mu^t}$, we extend this argument to consider $\ctype(P(w))^t) \rhd  \mu$.

\begin{proof}[Proof of  Theorem \ref{thm: desc basis are the same}]
We know $|\mathcal{C}_\mu| = |\mathcal{D}_{\mu^t}|$, since $|\mathcal{C}_\mu| = \text{dim}(R_{\mu})$ and $\mathcal{D}_{\mu^t}$ is a basis of $R_{\mu}$.
Hence to show equality of the two sets, we will show $\mathcal{C}_\mu \subset \mathcal{D}_{\mu^t}$.
From Corollary \ref{cor:ctype implies shuffle}, it suffices to consider $w\in S_n$ where $\ctype(P(w)^t) \rhd \mu$.
First, consider $w\in S_n$ such that $\ctype(P(w)^t)  = \lambda\unrhd \mu$ where $\lambda$ and $\mu$ differ by moving one box in the Young diagram. 
Assume that we construct $\mu$ from $\lambda$ by taking the last box in some column $j_2$ of $\lambda$ and moving it to the top of column $j_1$, where $j_1< j_2$. 


Since $\ctype(P(w)^t) = \ctype(\rev(w))$, we consider the filling $T_{\rev(w)}$ of $\lambda$ that we get from Algorithm \ref{alg: cat insertion}.
Let $a$ be the entry in the last box of column $j_2$. 
Take the box containing $a$ and append it into the end of column $j_1$.

By assumption, we know the resulting shape is $\mu$. Denote the resulting filling of shape $\mu$ by $T_{\rev(w)}^{(\mu)}$. 
We claim that each subword of $\cc(\rev(w))$ that corresponds to a column of $T_{\rev(w)}^{(\mu)}$ is still a valid cocharge word.
The only columns we need to consider are columns $j_1$ and $j_2$. Since the only modification to column $j_2$ is taking off the last entry we added, we see that the resulting word is still a valid cocharge word from Proposition \ref{prop: subword of column is a cocharge word}.

Since $j_1< j_2$ and $T_{\rev(w)}$ is partition shape, there exists an entry $b$ in column $j_1$ which is in the same row as $a$ in $T_{\rev(w)}$. Since $j_1<j_2$, we know we first added $b$, and then $a$, to the filling $T_{\rev(w)}$.  If we let $k_a, k_b$ denote the number of times we read $a,b$ before adding them to $T_w$, we must have $k_b \leq  k_a$. Thus, we know that $\cc(\rev(w))_a \leq  \cc(\rev(w))_b$ from Lemma \ref{lemma:row is sum of cocharge label and mult}. 

Thus adding $a$ to $\cc(\text{rev}(w))^{(j_1)}$ does not introduce a new cocharge value to this word since $\cc(\rev(w))_b$ is in $\cc(\text{rev}(w))^{(j_1)}$.
From case (i) of Lemma \ref{lemma: adding a repeated cocharge value is still a cc word}, we conclude the resulting word is a cocharge word of a permutation. 

If $\ctype(\rev(w))\unrhd \mu$ differ by moving multiple boxes, we repeat this process until we get a filling of $\mu$ where each of the columns correspond to valid cocharge words. It is clear that $\cc(\rev(w))$ is a shuffle of these cocharge words.
Hence $\cw(w) \in D_{\mu^t}$.
From this, we have $\mathcal{C}_\mu = \mathcal{D}_{\mu^t}$. Since $\mathbf{x}^{\mathcal{D}_{\mu^t}}$ is a basis of $R_{\mu}$, we conclude that our set  $\mathbf{x}^{\mathcal{C}_\mu}$  is a basis as well.
\end{proof}
